\chapter{La fuite}
\label{chap:la-fuite}

Khaer Ghal ne dormit pas.

Il resta longtemps assis contre la paroi de l'enclos, les jambes repliées
devant lui, sans savoir depuis combien de temps la nuit était tombée. À
quelques pas, Kamatsu n'avait presque pas bougé depuis qu'on les avait
enfermés.

Les deux garçons avaient peu parlé.

Khaer Ghal avait plusieurs fois cherché quelque chose à dire, mais chaque
phrase qui lui venait lui paraissait inutile. Il pouvait expliquer qu'il avait
voulu les aider, qu'il avait préparé leur fuite et qu'il n'avait jamais imaginé
qu'elle se terminerait ainsi. Rien de tout cela ne rendrait ses parents à
Kamatsu.

Alors il était resté silencieux.

Le village avait peu à peu sombré dans le même silence. Les dernières voix
s'étaient éteintes avec les feux et seuls les pas occasionnels d'un garde
rappelaient que tous les Orcs ne dormaient pas. Khaer Ghal avait fini par
reconnaître le rythme des rondes et suivait machinalement leurs passages,
écoutant les pas approcher, longer l'enclos puis s'éloigner entre les
habitations.

Il entendait justement l'un des gardes disparaître lorsqu'un autre bruit attira
son attention.

Des pas, plus légers cette fois.

Quelqu'un approchait.

Khaer Ghal se redressa et distingua bientôt une silhouette dans l'obscurité. Il
la reconnut avant même d'apercevoir son visage.

— Maman ?

Elle porta aussitôt un doigt à ses lèvres.

Khaer Ghal se tut.

Sa mère jeta un regard dans la direction prise par le garde, attendit encore
quelques secondes, puis sortit une clé de sous son vêtement. La serrure émit un
léger claquement qui parut beaucoup trop fort dans le silence de la nuit.

Kamatsu avait relevé la tête. Il observa la mère de Khaer Ghal, puis la porte
désormais ouverte, sans sembler comprendre ce qui se passait.

Elle entra dans l'enclos.

— Debout. Tous les deux.

Khaer Ghal se leva.

— Qu'est-ce que tu fais ?

— Je vous fais sortir.

— Mais père...

— Justement.

Elle se tourna vers lui et, même dans l'obscurité, Khaer Ghal comprit à son
expression qu'elle n'était pas venue discuter.

— Je ne sais pas ce qu'il a décidé de faire de toi au matin. Je sais seulement
ce que je l'ai vu faire depuis que tu as commencé à t'opposer à lui, et je ne
vais pas attendre de découvrir jusqu'où il est prêt à aller.

Khaer Ghal resta silencieux.

Kamatsu se leva.

— Vous venez avec nous ?

La mère de Khaer Ghal posa les yeux sur lui.

— Non.

Un seul mot, prononcé sans hésitation.

Puis elle ressortit de l'enclos.

— Venez.

\scenebreak

Ils traversèrent le village sans lumière.

La mère de Khaer Ghal avançait devant eux avec une assurance qui surprit son
fils. Elle connaissait les rondes aussi bien que lui, peut-être mieux encore,
et semblait savoir exactement quand traverser un espace découvert, quand
s'arrêter dans l'ombre d'une habitation et quand reprendre sa marche.

Khaer Ghal la suivait à quelques pas. Kamatsu venait derrière lui.

Personne ne parla.

À deux reprises, sa mère leva simplement une main et les trois silhouettes
s'immobilisèrent. Des guerriers passèrent au loin sans les voir. Ils
attendirent que leurs pas disparaissent avant de poursuivre leur chemin.

Peu à peu, les habitations se firent moins nombreuses et la haute palissade
apparut devant eux.

Khaer Ghal comprit aussitôt où sa mère les conduisait.

La brèche.

La portion qu'ils avaient démontée quelques heures auparavant n'avait pas
encore été réparée. Plusieurs rondins gisaient toujours au sol et l'ouverture
irrégulière découpait dans la fortification un passage assez large pour qu'un
adulte puisse le franchir rapidement en se courbant légèrement. Au-delà, la
forêt apparaissait entre les lourds troncs encore dressés de part et d'autre.

Mais quelque chose d'autre était resté.

Kamatsu s'arrêta.

Son père gisait encore près de l'ouverture.

Khaer Ghal entendit le souffle de son ami se couper. Kamatsu resta quelques
secondes parfaitement immobile, puis s'avança lentement vers le corps de
l'homme qui, quelques heures auparavant, avait choisi de défendre cette même
brèche pour donner à sa famille le temps de fuir.

Il s'agenouilla auprès de lui.

Khaer Ghal demeura en retrait. Sa mère aussi.

Il n'y avait rien à dire.

Kamatsu posa une main sur le bras de son père et resta ainsi un long moment, la
tête baissée. Lorsqu'il se releva enfin, son visage avait changé. Il essuya
rapidement ses yeux du revers de sa manche et regarda la forêt de l'autre côté
de la palissade.

La mère de Khaer Ghal attendit encore quelques secondes.

Puis elle posa doucement une main dans son dos.

— Il faut y aller.

Kamatsu acquiesça.

Cette fois, ils franchirent la brèche.

\scenebreak

De l'autre côté de la palissade, la forêt commençait presque immédiatement. La
lumière froide de la lune descendait entre les branches tandis que, derrière
eux, quelques lueurs orangées du village filtraient encore à travers
l'ouverture.

La mère de Khaer Ghal leur remit d'abord ce qu'elle avait préparé. Un petit sac
contenait de la nourriture, une outre, une couverture et quelques affaires
qu'elle avait pu rassembler sans attirer l'attention. Ce n'était pas
grand-chose, mais suffisamment pour leur permettre de tenir un jour ou deux
s'ils se montraient prudents.

Elle tendit ensuite une épée courte à son fils.

Khaer Ghal la reconnut immédiatement.

C'était celle avec laquelle il s'était entraîné pendant une partie de son
apprentissage, une arme orque simple et robuste dont il connaissait déjà le
poids, l'équilibre et jusqu'aux irrégularités de la poignée. Il passa le
fourreau à sa ceinture et leva les yeux vers sa mère.

Elle avait encore quelque chose à lui donner.

Ses doigts se refermèrent autour du cordon qu'elle portait depuis aussi
longtemps qu'il pouvait s'en souvenir. Une petite pierre sombre y était
suspendue, marquée d'un ancien symbole orc dont les lignes irrégulières avaient
été adoucies par les années.

Khaer Ghal l'avait vue des centaines de fois.

Il ne l'avait jamais vue quitter le cou de sa mère.

— Maman...

Elle passa le cordon par-dessus sa tête et laissa la pierre retomber contre sa
poitrine.

Khaer Ghal baissa les yeux. Ses doigts se refermèrent instinctivement autour du
pendentif, suffisamment fort pour sentir ses contours s'enfoncer dans sa paume.

Lorsqu'il releva la tête, sa mère avait posé ses deux mains sur ses épaules.

Ils se regardèrent.

Pendant quelques secondes, le village, la fuite et le danger semblèrent
s'éloigner. Khaer Ghal ne voyait plus que le visage de sa mère. Il voulait lui
demander de venir. Il voulait lui demander ce qui lui arriverait lorsqu'on
découvrirait leur disparition, ce qu'elle dirait à son père et s'ils se
reverraient.

Aucune de ces questions ne franchit ses lèvres.

Sa mère resserra légèrement ses mains sur ses épaules.

— Écoute-moi bien, Khaer Ghal.

Il acquiesça.

— Ne laisse jamais personne choisir ta route à ta place.

Khaer Ghal resta silencieux.

Il était encore trop jeune pour mesurer tout ce que ces mots finiraient par
représenter. Quelques jours auparavant, il n'aurait peut-être même pas compris
ce qu'elle voulait lui dire. Toute sa vie, quelqu'un avait décidé de ce qu'il
devait apprendre, de ce qu'il devait devenir, de ceux auxquels il pouvait
parler et de la place qu'il devait occuper.

Puis Kamatsu lui avait demandé :

\emph{Et si on ne veut pas ?}

Khaer Ghal avait commencé à comprendre.

Cette nuit-là, sa mère lui donnait une réponse.

Elle leva une main jusqu'à sa joue.

— Tu feras des erreurs. Tu prendras de mauvaises décisions. Tu tomberas
peut-être.

Un sourire triste passa sur son visage.

— Mais qu'elles soient au moins les tiennes.

Khaer Ghal sentit ses yeux le brûler. Une larme finit par rouler sur sa joue,
mais il ne détourna pas le regard.

— Viens avec nous.

Sa mère secoua doucement la tête.

— Je ne peux pas.

— Pourquoi ?

Elle regarda un instant la palissade derrière eux.

— Parce que ma vie est ici.

Khaer Ghal voulut protester, mais elle l'attira contre elle avant qu'il puisse
répondre.

Cette fois, il la serra immédiatement dans ses bras.

— Je t'aimerai toujours, murmura-t-elle.

Khaer Ghal ferma les yeux.

Il aurait voulu que cet instant dure davantage.

Ce fut pourtant elle qui finit par se détacher. Ses mains retrouvèrent les
épaules de son fils et elle le regarda encore une fois droit dans les yeux.

Khaer Ghal gardait toujours une main refermée sur le pendentif.

Un peu en retrait, Kamatsu les observait en silence.

\illustration
  {0700.png}
  {Les adieux}
  {la-fuite}

La mère de Khaer Ghal tourna légèrement la tête vers lui.

— Veille sur ton ami.

Khaer Ghal suivit son regard jusqu'à Kamatsu.

Le mot le surprit.

Il n'avait jamais réellement cherché à donner un nom à ce qui s'était installé
entre eux au fil des derniers jours. Pourtant, entendu ainsi, après l'arène,
leur tentative de fuite et cette nuit passée enfermés ensemble, il ne lui parut
pas faux.

Sa mère regarda alors Kamatsu.

— Et toi, protège-le.

Khaer Ghal releva les yeux, surpris.

Depuis le début, il lui avait semblé évident que c'était à lui de protéger
Kamatsu. Il était l'Orc. Il connaissait la forêt, savait combattre et avait déjà
essayé de sauver son ami à deux reprises.

Sa mère venait pourtant de confier à Kamatsu exactement la même responsabilité.

Le garçon regarda Khaer Ghal, puis sa mère.

— Je le ferai.

La réponse fut simple.

Sa mère hocha la tête.

Puis elle recula d'un pas.

Le moment était terminé.

— Maintenant, partez.

\scenebreak

Les deux garçons s'enfoncèrent dans la forêt.

Khaer Ghal connaissait suffisamment les environs pour choisir une direction et
éviter les sentiers les plus évidents. Pourtant, jamais il ne s'était éloigné
du village de cette manière. Lorsqu'il était parti chasser, explorer ou
s'entraîner dans les bois, il avait toujours existé quelque part derrière lui
un endroit où revenir. Quelqu'un savait qu'il était parti. Quelqu'un attendait
son retour.

Cette fois, rien ne l'attendait.

La sensation aurait dû être effrayante.

Elle l'était.

Mais pas seulement.

Après un long moment de marche, Khaer Ghal s'arrêta. Kamatsu fit encore
quelques pas avant de remarquer qu'il ne le suivait plus et revint près de lui.

Tous deux se retournèrent.

À travers les arbres, on distinguait encore par endroits la palissade et
quelques lueurs provenant du village. La brèche n'était plus qu'une tache
sombre entre les troncs.

Sa mère avait disparu.

Khaer Ghal écouta attentivement.

Aucun cri.

Aucun appel aux armes.

Aucune agitation.

Rien n'indiquait encore que leur fuite avait été découverte.

Il ne savait pas si sa mère avait réussi à regagner sa couche sans être vue. Il
ignorait ce qu'elle raconterait au matin et ce que son père ferait lorsqu'il
découvrirait leur disparition. Plus encore, il ne savait pas ce qu'il ferait
lui-même le lendemain, ni où ils dormiraient la nuit suivante.

Jusqu'à présent, il avait toujours connu la réponse à ces questions.

Cette nuit, pour la première fois, personne ne la connaissait à sa place.

Khaer Ghal baissa les yeux vers la pierre suspendue à son cou et repensa aux
paroles de sa mère.

\emph{Ne laisse jamais personne choisir ta route à ta place.}

Il referma les doigts autour du pendentif.

Puis il se détourna du village.

Kamatsu regardait déjà la forêt devant eux.

— On va où ? demanda-t-il.

Khaer Ghal observa les arbres, les reliefs sombres qu'il distinguait entre les
troncs et le ciel nocturne au-dessus des branches.

Il connaissait les environs.

Au-delà, beaucoup moins.

— Par là.
